% =============================================================================
%  Abstract (arXiv: <= 1920 символов в форме подачи; текст в PDF может совпадать)
%  Структура: контекст -> проблема -> что сделано -> ключевые числа -> вывод
% =============================================================================
\begin{abstract}
Simultaneous Localization and Mapping (SLAM) is a core capability of
autonomous robots, drones, and AR/VR headsets, yet state-of-the-art
algorithms remain too power-hungry for milliwatt-class edge platforms. The bottleneck of modern Edge AI SoC is memory bus and DMA are highly loads while data transferring before computation, also computation blocs such as NPU/TPU/GPU have general architecture and not optimized for SLAM. Specifically to SLAM and other sensor data processing, moving data from camera to CPU, feature extraction and then passing to computation block dominate latency and energy on SoC. In this paper we present \sysname, a \TODO{process node, e.g., 28-nm CMOS}
ASIC accelerator for \TODO{visual / visual-inertial} SLAM that
\TODO{key idea 1, e.g., a streaming feature front-end} and
\TODO{key idea 2, e.g., a sparse Schur-complement back-end with on-chip
marginalization}.
The chip occupies \NUM~\si{\milli\metre\squared}, runs at \NUM~\si{\mega\hertz},
and processes \NUM\fps\ VGA input while consuming \NUM~\si{\milli\watt}
--- \NUM$\times$ more energy-efficient than \TODO{baseline platform}
with \TODO{no / $<$X\%} accuracy degradation on the EuRoC and TUM RGB-D
benchmarks.
\TODO{One-sentence take-away / availability of RTL or dataset.}
\end{abstract}

\noindent\textbf{Keywords:} SLAM, ASIC, hardware accelerator, visual odometry,
bundle adjustment, energy efficiency, edge robotics.
